{\centering
	\par
	\large \textbf{\underline{ABSTRACT}} \par
}\vspace{0.7cm}

\noindent Modern cybersecurity landscapes face growing challenges from sophisticated kernel-level rootkits, including BPFDoor, Syslogk, cd00r, Reptile, and Jynx, which leverage covert network signaling mechanisms commonly referred to as ``magic packets.'' These specially constructed packets carry concealed trigger sequences within TCP header fields, payload byte patterns, and protocol-level anomalies, enabling adversaries to silently activate hidden backdoors and bypass conventional security tools such as signature-based Intrusion Detection Systems and log-monitoring frameworks.\\

\noindent To address these challenges, this work introduces \textbf{Sentinel Forge}, an intelligent and self-operating network defense platform that leverages multi-layer Deep Packet Inspection to identify, classify, and counteract rootkit-driven covert traffic in near real time. The platform adopts an Out-of-Band monitoring design combined with a centralized enforcement module termed the Universal Agent Core, ensuring that production network performance remains entirely unaffected during analysis while preserving comprehensive visibility across OSI Layers 3, 4, and 7.\\

\noindent A key contribution of this work is the development of an integrated anti-evasion processing pipeline that performs IP fragment reconstruction, Base64 and hexadecimal payload decoding, exhaustive single-byte XOR key recovery, and Shannon Entropy computation to uncover encrypted or obfuscated content. Threat classification is performed through a feature-weighted risk-scoring model, while a trained Decision Tree classifier serves as a rapid pre-screening filter. Secure coordination between detection and response components is achieved via HMAC-SHA256 signed webhooks, complemented by LRU-based rate throttling and a modular firewall plugin system compatible with pfSense and Linux iptables.\\

\noindent Testing across seven rootkit families yielded a detection rate surpassing 90\%, zero false positives against benign traffic baselines, and a mean detection-to-response time below 50 milliseconds over 100 trials. The findings confirm meaningful architectural improvements over both Fail2Ban and conventional inline firewall solutions.\\

\noindent\textbf{Keywords:} Deep Packet Inspection, Rootkit Detection, Magic Packet, Shannon Entropy, Machine Learning, Intrusion Detection System, Out-of-Band Architecture, Network Security, Anti-Evasion, Automated Threat Mitigation